\section{Strategic information and independent safeguards}\label{app:strategy}
An audit plan specifies which physical command to invoke at each observable record and when to stop. A fixed plan alone may reveal nothing about one pair of worlds even when privately mixing two plans yields a useful test. Conversely, distinguishing each pair with some test can still leave mixtures of the two classes indistinguishable.

\begin{theorem}[A.G1. The finite complete-record minimax value]\label{adv:ag1}
Let \(S\) be the union of the finite stopped records of every legal deterministic complete acquisition-and-stop plan \(d\) in a nonempty finite class \(\mathcal D\). Every stop permits either terminal label. Include clocks, prices, delivered signals, any public announcements, commands and legal startup. For each fixed world \(w\), define \(x_w(h)\) as the probability of the stopped record's data under its prescribed sequence of commands. For \(h\in S_d\), the stopped records of \(d\), these probabilities sum to one. Source and responder depend only on actually available inputs; future private plans and hidden seeds are not observable. For nonempty target classes \(b\in\{+,-\}\) put \(K_b=\overline{\operatorname{conv}}\{x_w:Y(w)=b\}\).
With arbitrary private design and terminal-report randomization, the minimum worst-world binary error is attained and equals
\[
 \frac{1-\gamma}{2},\qquad
 \gamma=\min_{\substack{0\le\lambda\le1\\p\in K_+,\,q\in K_-}}
       \max_{d\in\mathcal D}\sum_{h\in S_d}
                   |\lambda p(h)-(1-\lambda)q(h)|.
\]
In particular an error \(\delta\le1/2\) is possible iff \(\gamma\ge1-2\delta\). Chance performance occurs iff the hulls intersect. With compact class images, each mixture witness needs at most \(|S|+1\) realizable worlds per label; otherwise every strict lower error bound below the minimax error has such finite witnesses, without necessarily attaining equality. For finite rational worlds and a factored local legal menu, a rational sequence program is polynomial in the expanded tree and world data.
\end{theorem}
\begin{proof}
A pure policy is a complete design and a terminal labeling. A draw of the private seed fixes both tables in advance; admissible mixtures over such tables are implementable. Fix a prior mass \(\lambda\) and a law in each hull. The minimum error for design \(d\) is the sum on \(S_d\) of the smaller of the two weighted masses. Their total sum is one, so the risk is one half of one minus the displayed norm. Minimize over designs and maximize over adversarial mixtures: the finite-dimensional minimax theorem applies to the compact policy simplex and the compact convex set of weighted law pairs. Closing and convexifying neither increases nor decreases the worst linear payoff of a fixed policy. The extrema are attained there. Zero norm forces equality at every relevant coordinate; normalization at any design forces \(\lambda=1/2\). Carathéodory's theorem gives the support bound in compact images. Approximate a least favorable closed-hull law by finite convex combinations for a strict margin otherwise. The finitely many design losses are continuous, so the strict margin survives. The rational finite program uses Theorem~\ref{flag:u1} on the legal pre-verdict tree. Global constraints on complete plans may instead require their explicit list.
\end{proof}
To see the two effects separately, let the probabilities of a bit being one under probes \(A,B\) be \((1/4,3/4),(3/4,1/4)\) on positives and \((1/8,1/8),(7/8,7/8)\) on negatives. Equal mixtures agree even though each opposing pair differs on both probes. Alternatively use a single positive \((1/2,1/2)\) and negatives \((1/2,1),(1,1/2)\). Positive prior \(3/5\) versus equal negative mixtures gives Bayes error \(2/5\) on each probe. Choose probes fairly, report positive after zero, and negative with probability \(4/5\) after one. Each error is \(2/5\). No single predetermined probe improves on chance in the worst case.

\begin{theorem}[A.G2. Independent facts that restore a feasible policy]\label{adv:ag2}
Fix \(m\ge1\) authorized deterministic complete policies and their private mixtures in the simplex \(\Delta_m\). Each world's closed convex constraint set \(P_w\subseteq\Delta_m\) is defined by affine performance requirements on genuine quantities. Supported-path prohibitions can be zero-probability tests where support is exact; check hard permission separately where that encoding is insufficient. Each nonchanging verified fact has a satisfaction set \(F_j\) and additive price \(c_j\ge0\). Both designated, physically realizable opposite-label reference worlds satisfy all the facts, and a positive service floor holds as an explicit requirement at the positive reference. A selected fact set \(J\) admits one feasible common policy exactly when it contains a fact excluding a member of every group \(H\) of at most \(m\) original worlds with \(\bigcap_{w\in H}P_w=\varnothing\). Minimum realized price is the corresponding weighted hitting-set minimum. Fact selection is NP-hard even with two-world minimal obstructions.
\end{theorem}
\begin{proof}
A surviving infeasible group rules out a common point. If the intersection over all survivors is empty, compactness of the simplex yields a finite witness and Helly gives at most \(m\) worlds in affine dimension \(m-1\). Fixed finite rational data permit exact intersection tests and enumeration or separation of blockers. To reduce set cover, keep a clean positive that returns bit one and can deliver a safe job, a clean negative returning zero, and one physically faulty negative per universe element which returns bit one and whose finish violates the event. Require correct decisions, zero risk and positive useful service in the clean positive. Each faulty world blocks jointly with the clean positive. Let a fact exclude exactly the faulty worlds in its subset by checking the responsible maintenance routes. The clean references satisfy every fact. A feasible selection is exactly a set cover, with the same price. The selected facts leave the surviving worlds' physical rows and controller class intact.
\end{proof}

\begin{theorem}[A.G3. Different instruments on different routes]\label{adv:ag3}
Let an installation fix a nominal type \(b\in\{+,-\}\), \(s\in[0,1]\), and a linked response rule. Safe intake delivers a fair left/right token, independent of the world. On left, one diagnostic costs one; on right the only diagnostic costs three. With either, the responder chooses \(j\in\{0,1\}\) before a fresh draw. Its probability of reporting the true type is \(3/4+(s+k_b j)/32\) on left and \(7/8+(s+k_b j)/64\) on right, where \(k_+=1,k_-=2\). No extra pre-report information is available. The controller may also skip and may abstain. A positive declaration triggers the actual routed finish once. After seeing that verdict the same responder chooses \(a\in\{0,1\}\) before a fresh physical outcome. Conditional defect probability is the entry below plus \((s+a)/16\):
\[
\begin{array}{c|cc}\text{type}&L&R\\ \hline
+&1/8&1/4\\ -&5/8&3/4
\end{array}
\]
The intrinsic label is the direct always-finishing reference risk at cutoff \(1/2\). Require declared-wrong probability at most \(5/32\), declaration probability at least \(3/4\) and whole physical risk at most \(1/4\). On positives also require issuance at least \(19/32\), selective risk at most \(3/8\) and expected safe routed usefulness at least \(95/256\). With the complete record, exact laws and matched \(s=j=0\) worlds, the optimal worst expected acquisition is \(5/4\) and the minimum hard price is three. Precommitting to one probe identity before receiving the token is infeasible.
\end{theorem}
\begin{proof}
The direct risk on the two types ranges respectively in \([3/16,5/16]\) and \([11/16,13/16]\), so the target is \(b\). On an equal prior on the two matched worlds let \(\alpha,\beta\) denote the conditional purchase probabilities on left and right. A declared verdict without information gains zero over a fair guess; a purchase can gain at most \(1/4\) or \(3/8\). Denote by \(D,E\) a declaration and a wrong declaration. Thus
\[
 7/32\le\Pr(D)/2-\Pr(E)\le\alpha/8+3\beta/16.
\]
Consequently \(4\alpha+6\beta\ge7\), \(\beta\ge1/2\), and \((\alpha+3\beta)/2\ge5/4\). A precommitted identity yields \(\alpha+\beta\le1\), making the inequality impossible. Buy on every left and half the rights, declare from the bought bit, otherwise abstain. This attains both prices and \(3/4\) declaration. Each positive cut has defect at most \(3/8\); wrong-declaration probability is at most \((1/2)(1/4)+(1/4)(1/8)=5/32\). Positive issuance exceeds or equals \(19/32\) and safe service is at least \((19/32)(5/8)\). Whole risk is at most \(7/32\) on positives, \(35/256\) on negatives. On the matched negative, each issuance has conditional defect at least \(5/8\) and shares its full pre-report support with the matched positive. Uniform selective risk \(1/2\) across both signs cannot coexist with positive exposure.

For the realized safeguard example fix \(s=0\) and add three possible maintenance branches which make every finish defective without changing pre-report data. Returning bit-identical records in a clean positive and a corresponding faulty negative forces their two wrong-label probabilities to sum to the common declaration probability, at least \(3/4>2(5/32)\). The three available facts exclude branch pairs \(\{1,2\},\{2,3\},\{1,3\}\) at prices \(2,3,4\). All are true in the clean installation. Each one leaves a blocker; the cheapest two remove all of them and retain both clean signs. Their realized price five and the per-attempt probe price answer separate questions.
\end{proof}
